% ============================================================
% Exponential probability graph
% ============================================================

\newcommand{\ExponentialProbabilityGraph}[3]{%
	% #1 = lambda
	% #2 = lower bound a
	% #3 = upper bound b
	
	\begin{center}
		\begin{tikzpicture}[x=1cm,y=10cm]
			
			% Graph limits
			\pgfmathsetmacro{\GraphXMax}{#3+1.2}
			
			% Maximum value of the exponential density
			\pgfmathsetmacro{\GraphYMax}{1.2*#1}
			
			% Function values
			\pgfmathsetmacro{\Fa}{#1*exp(-#1*#2)}
			\pgfmathsetmacro{\Fb}{#1*exp(-#1*#3)}
			
			% Axes
			\draw[->] (-0.5,0) -- (\GraphXMax,0) node[right] {$t$};
			\draw[->] (0,-0.02) -- (0,\GraphYMax);
			
			% Shaded probability region
			\fill[blue!25, fill opacity=0.7]
			(#2,0)
			-- plot[
			domain=#2:#3,
			samples=100
			]
			(\x,{#1*exp(-#1*\x)})
			-- (#3,0)
			-- cycle;
			
			% Exponential density
			\draw[
			thick,
			domain=0:\GraphXMax,
			samples=100
			]
			plot (\x,{#1*exp(-#1*\x)});
			
			% Vertical guides
			\draw[dashed]
			(#2,0) -- (#2,\Fa);
			
			\draw[dashed]
			(#3,0) -- (#3,\Fb);
			
			% Horizontal axis labels
			\node[below=2pt] at (#2,0) {\scriptsize $#2$};
			\node[below=2pt] at (#3,0) {\scriptsize $#3$};
			
			% Probability label
			\pgfmathsetmacro{\MidPoint}{(#2+#3)/2}
			\node at (\MidPoint,0.10)
			{\scriptsize $\displaystyle P(#2<T<#3)$};
			
			% Density label
			\pgfmathsetmacro{\LabelX}{max(1.5,min(#3-0.3,3.5))}
			\pgfmathsetmacro{\LabelY}{#1*exp(-#1*\LabelX)+0.06}
			
			\node[above] at (\LabelX,\LabelY)
			{\scriptsize $\displaystyle f_T(t)=#1e^{-#1t}$};
			
		\end{tikzpicture}
	\end{center}
}

% ============================================================
% Exponential probability graph: less than a
% ============================================================

\newcommand{\ExponentialLessThanGraph}[3]{%
	% #1 = lambda
	% #2 = upper bound a
	% #3 = graph extension beyond a
	
	\begin{center}
		\begin{tikzpicture}[x=1cm,y=10cm]
			
			% Graph limits
			\pgfmathsetmacro{\GraphXMax}{#2+#3}
			
			% Maximum value of the exponential density
			\pgfmathsetmacro{\GraphYMax}{1.2*#1}
			
			% Function value at the cutoff
			\pgfmathsetmacro{\Fa}{#1*exp(-#1*#2)}
			
			% Axes
			\draw[->] (-0.5,0) -- (\GraphXMax,0) node[right] {$t$};
			\draw[->] (0,-0.02) -- (0,\GraphYMax);
			
			% Shaded probability region
			\fill[blue!25, fill opacity=0.7]
			(0,0)
			-- plot[
			domain=0:#2,
			samples=100
			]
			(\x,{#1*exp(-#1*\x)})
			-- (#2,0)
			-- cycle;
			
			% Exponential density
			\draw[
			thick,
			domain=0:\GraphXMax,
			samples=100
			]
			plot (\x,{#1*exp(-#1*\x)});
			
			% Vertical guide at the cutoff
			\draw[dashed]
			(#2,0) -- (#2,\Fa);
			
			% Horizontal axis label
			\node[below=2pt] at (#2,0)
			{\scriptsize $#2$};
			
			% Probability label
			\pgfmathsetmacro{\LabelX}{0.45*#2}
			\pgfmathsetmacro{\LabelY}{0.10}
			
			\node at (\LabelX,\LabelY)
			{\scriptsize $\displaystyle P(T<#2)$};
			
			% Density label
			\pgfmathsetmacro{\DensityLabelX}{max(1.2,min(#2-0.5,3.5))}
			\pgfmathsetmacro{\DensityLabelY}
			{#1*exp(-#1*\DensityLabelX)+0.06}
			
			\node[above] at (\DensityLabelX,\DensityLabelY)
			{\scriptsize $\displaystyle f_T(t)=#1e^{-#1t}$};
			
		\end{tikzpicture}
	\end{center}
}

% ============================================================
% Exponential probability graph: greater than a
% ============================================================

\newcommand{\ExponentialGreaterThanGraph}[3]{%
	% #1 = lambda
	% #2 = lower bound a
	% #3 = graph extension beyond a
	%
	% Example:
	% \ExponentialGreaterThanGraph{0.20}{7}{5}
	%
	% Represents P(T > 7), with the graph extending
	% 5 units beyond the cutoff.
	
	\begin{center}
		\begin{tikzpicture}[x=1cm,y=10cm]
			
			% Graph limits
			\pgfmathsetmacro{\GraphXMax}{#2+#3}
			
			% Maximum value of the exponential density
			\pgfmathsetmacro{\GraphYMax}{1.2*#1}
			
			% Function value at the cutoff
			\pgfmathsetmacro{\Fa}{#1*exp(-#1*#2)}
			
			% Axes
			\draw[->] (-0.5,0) -- (\GraphXMax,0) node[right] {$t$};
			\draw[->] (0,-0.02) -- (0,\GraphYMax);
			
			% Shaded probability region
			\fill[blue!25, fill opacity=0.7]
			(#2,0)
			-- plot[
			domain=#2:\GraphXMax,
			samples=100
			]
			(\x,{#1*exp(-#1*\x)})
			-- (\GraphXMax,0)
			-- cycle;
			
			% Exponential density
			\draw[
			thick,
			domain=0:\GraphXMax,
			samples=100
			]
			plot (\x,{#1*exp(-#1*\x)});
			
			% Vertical guide at the cutoff
			\draw[dashed]
			(#2,0) -- (#2,\Fa);
			
			% Horizontal axis label
			\node[below=2pt] at (#2,0)
			{\scriptsize $#2$};
			
			% Probability label
			\pgfmathsetmacro{\LabelX}{#2+0.35*#3}
			\pgfmathsetmacro{\LabelY}{0.10}
			
			\node at (\LabelX,\LabelY)
			{\scriptsize $\displaystyle P(T>#2)$};
			
			% Density label
			\pgfmathsetmacro{\DensityLabelX}{max(1.2,min(#2-0.5,3.5))}
			\pgfmathsetmacro{\DensityLabelY}
			{#1*exp(-#1*\DensityLabelX)+0.06}
			
			\node[above] at (\DensityLabelX,\DensityLabelY)
			{\scriptsize $\displaystyle f_T(t)=#1e^{-#1t}$};
			
		\end{tikzpicture}
	\end{center}
}